\chapter{Experiment Design}\label{ch:design}

\section{Introduction}
The experiments of this report follow one design, \emph{setup~v2} (fixed on 22~November 2026): bench experiments that
compare \drtc{} with LeRobot's client under every fault type, take \drtc{} apart and choose a smoothing stack, and a
robot experiment that tests \drtc{} and the stack on the real arm. Every trained policy also has a baseline on the
robot. This chapter describes the design, the claims it tests and the protocol.
\textslot{one paragraph on the earlier pilot run (20~November) and why it led to setup~v2, without its results.}

\section{Setup v2}\label{sec:setupv2}

\begin{figure}[htbp]
  \centering
  % Setup v2: the bench experiments feed the choice of the stack, which the robot experiment uses.
% Shared figure colours (colour only in figures).
\definecolor{robotfill}{HTML}{EAF2FC}\definecolor{robotline}{HTML}{2A78D6}
\definecolor{serverfill}{HTML}{FDEEE7}\definecolor{serverline}{HTML}{EB6834}
\definecolor{netfill}{HTML}{F0EFEC}\definecolor{greenfill}{HTML}{E9F5EC}\definecolor{greenline}{HTML}{2F9E58}
\definecolor{purplefill}{HTML}{F1ECF9}\definecolor{purpleline}{HTML}{7A55C7}
\definecolor{ink}{HTML}{0B0B0B}\definecolor{ink2}{HTML}{52514E}\definecolor{faultc}{HTML}{C9C6BF}

\begin{tikzpicture}[font=\small, node distance=5mm and 8mm,
  b/.style={draw=#1, thick, rounded corners=2pt, fill=white, align=left, text width=40mm, inner sep=4pt},
  arr/.style={-{Latex[length=2.2mm]}, thick, ink2}]
\node[b=greenline] (e12) {\textbf{E1 + E2} (bench)\\ upstream patient, as shipped, \drtc;\\ \drtc{} without each part\\ $\times$ 7 fault types};
\node[b=greenline, below=of e12] (e3a) {\textbf{E3a} (bench)\\ each smoothing method alone\\ on \drtc{} + mix};
\node[b=purpleline, right=of e3a] (rule) {\textbf{The stack rule}\\ lowest jerk, unless\\ M12 +1\degree{} or\\ M2 / M11 +10\%};
\node[b=greenline, right=of rule, yshift=14mm] (e3b) {\textbf{E3b} (bench)\\ the stack, and the stack\\ without each part};
\node[b=robotline, right=of rule, yshift=-14mm] (e4) {\textbf{E4} (robot, success)\\ patient upstream, + mix,\\ \drtc{} + mix, + stack + mix};
\draw[arr] (e3a) -- (rule); \draw[arr] (rule.east) -- ++(3mm,0) |- (e3b.west); \draw[arr] (rule.east) -- ++(3mm,0) |- (e4.west);
\node[right=8mm of e12, text width=60mm, font=\scriptsize, text=ink2, align=left] {Bench: replay robot, simulated network\\ sorting + cube in box v1, five policies, 8 runs per cell};
\end{tikzpicture}

  \caption{Setup v2. The bench experiments E1--E3a run first; E3a chooses the smoothing stack by a rule fixed in
           advance; E3b tests the stack on the bench and E4 on the robot.}
  \label{fig:design}
\end{figure}

\paragraph{The bench} runs on the sorting and the cube-in-box v1 tasks for all seven policies (\smolvla, \groot,
GR00T~N1.7, ACT, Diffusion Policy, \pizerofive{} and MolmoAct~2), eight runs per cell, 14~s each.
\begin{itemize}
\item \textbf{E1, \drtc{} against upstream:} upstream patient, upstream as shipped and \drtc, under every fault type of
      Table~\ref{tab:faults} (as shipped without the fault-free cell).
\item \textbf{E2, ablations:} \drtc{} without its cooldown (A1), with upstream's latency estimate (A2), with upstream's
      trigger (A3), with the original $H/2$ cooldown cap (A4, which undoes our first fix), under every fault type; and
      $s_{\min}$ at half and one and a half times its value (A5), with no fault and under the mix.
\item \textbf{E3a, smoothing methods:} \drtc{} under the mix with each method of Table~\ref{tab:smoothing} alone, and
      with the Butterworth filter plus adaptive playback (the first run's condition~4); RTC only for \smolvla{} and
      \groot.
\item \textbf{E3b, the stack:} the stack chosen by the rule of Section~\ref{sec:stackrule}, without faults and under the
      mix, and the stack without each of its parts. When the stack uses RTC, the policies without it are run with
      ensembling instead, in a plan of their own.
\end{itemize}
E1 and E2 share one plan per task (52 conditions), so every ablation is compared with \drtc{} inside the same sessions.

\paragraph{The robot experiment E4} runs the four policies with the best robot baselines (GR00T~N1.7, \smolvla,
Diffusion and \pizerofive, Section~\ref{sec:baselines}) on both tasks, ten trials per condition (160 trials per task),
on the real network: upstream patient without faults and under the mix, \drtc{} under the mix, and \drtc{} with the
stack under the mix. The no-fault condition is kept inside E4, although the baselines ran the same condition, so that
every comparison stays within one interleaved plan. The stack on the robot is the stack without RTC (adaptive playback,
temporal ensembling, the Butterworth filter and the resume ramp) for all four policies: Diffusion and GR00T~N1.7 cannot
use RTC, and one stack for every policy keeps the comparison simple; the choice is discussed with the bench results
(Section~\ref{sec:res-e3}). The stack with RTC is kept for a later run.

\figslot{fig:e4ladder}{The four conditions of E4 and what each comparison shows.}{D8.1 Diagram}{A ladder of the four
conditions: 1 patient LeRobot, normal network $\rightarrow$ 2 the same under the fault mix (what the faults cost)
$\rightarrow$ 3 \drtc{} under the mix (what \drtc{} wins back, the main result) $\rightarrow$ 4 \drtc{} with the stack
under the mix (what smoothing adds); with the fault mix drawn once (a 1~s latency spike, 1~s of lost replies, a 1.5~s
outage, shifted by 0 or 4~s).}

\paragraph{Claims and primary comparisons.} Table~\ref{tab:claims} lists the comparisons that are tested; everything
else is exploratory (Section~\ref{sec:stats}).

\begin{table}[htbp]
  \centering
  \small
  \begin{tabularx}{\textwidth}{X l l}
    \toprule
    Claim & Comparison & Metric \\
    \midrule
    Faults hurt even a client that waits & E4: upstream patient + mix against no fault & M1 \\
    \drtc's scheduling beats waiting & E4: \drtc{} + mix against upstream patient + mix & M1 \\
     & E1: \drtc{} against upstream patient, each fault type & M2, M11 \\
    Smoothing improves the motion without lowering success & E4: \drtc{} + stack against \drtc, under the mix & M1 \\
    \bottomrule
  \end{tabularx}
  \caption{The claims and their primary comparisons.}\label{tab:claims}
\end{table}

\section{Baselines}\label{sec:baselines-design}

Every trained policy is also evaluated on the robot on both tasks without faults, with patient upstream (the same
reference condition as E4), ten trials each, all seven trained policies in one plan per task; the cube-in-box v0
baseline of the first two policies is optional.

\section{Protocol}\label{sec:protocol}

\begin{itemize}
\item \textbf{Machines.} One robot computer and one model server serve a whole plan where possible: the robot
      baselines ran from the Raspberry Pi~5 with every model on the H200; every trial records its robot computer and
      server, and the timing metrics are never averaged across them.
\item \textbf{Sessions.} One session at a time, from one robot computer; each session pulls the results store first
      (or runs offline, as the Pi did over its slow WiFi) and pushes at the end. For E4 the Pi is connected to the
      laptop by an Ethernet cable: over the WiFi hotspot each request (two uncompressed camera images, about 2.9~MB)
      took about a second to send, and the replies of the slower policies came back after about 2~s, longer than any
      bench condition and enough for \drtc's cooldown to abandon every reply.
\item \textbf{Layouts.} One layout per round, rotated over the rounds as in the recording.
\item \textbf{Warm-up.} Plans with an RTC condition begin with one uncounted warm-up round (Section~\ref{sec:rtcmethod}).
\item \textbf{Growth.} When there is time, in this order: E4 to 20 trials; E4 with the RTC stack for \smolvla{} and
      \pizerofive; the baselines to 20 trials; robot runs of a bench fault type that behaved unexpectedly.
\end{itemize}

\section{Summary}
\textslot{the chapter's summary: setup v2, E4, the baselines, the claims and the protocol}
